% !TeX program = uplatex
% !TeX encoding = UTF-8 Unicode
\documentclass[a4paper,10pt]{jsarticle}
\usepackage[dvipdfmx]{graphicx}
\usepackage{otf}
\usepackage{amssymb}
\usepackage{amsmath}
\usepackage{float}
% \usepackage[compact]{titlesec}
\usepackage{fancyhdr}
\usepackage{url}

%---------------------------------------------------
% ページの設定
%---------------------------------------------------
\setlength{\textwidth}{179truemm}
\setlength{\textheight}{260truemm}
\setlength{\topmargin}{-14.5truemm}
\setlength{\oddsidemargin}{-9.5truemm}
\pagestyle{fancy}
\fancyhf{}
\fancyhead[L]{}
\fancyhead[C]{}
\fancyhead[R]{}
\renewcommand{\headrulewidth}{0pt}
\fancyfoot[L]{}
\fancyfoot[C]{}
\fancyfoot[R]{}
\renewcommand{\footrulewidth}{0pt}
\setlength{\headheight}{13pt}
\setlength{\parindent}{1zw}

\begin{document}

\begin{center}
  % 修士研究での発展性を感じさせるタイトル
  {\LARGE 曖昧移動指示に対応するナビゲーションシステム}
\end{center}
\vspace{5truemm}

\section{手法}

本研究では，HOYO 歩容データセットと日本語オノマトペから，人間の歩行の「質感」
（速い／遅い／重い／ふらふら）を表現する運動潜在表現を学習する．
学習された潜在表現は，Navila ベースの強化学習エージェントに対して
スタイル条件付きの報酬として用いることを想定している．

\subsection{データセットと前処理}

\paragraph{HOYO Instruction Dataset}
HOYO データセットから，日本語オノマトペ 11 種類に対応するサンプルのみを抽出し，
簡易データセット \texttt{HoyoInstructionDataset} を構成する．
ラベル集合は
\[
\{\text{通常},\ \text{すたすた},\ \text{せかせか},\ \text{てくてく},\ 
\text{どっしどっし},\ \text{とぼとぼ},\ \text{のしのし},\ \text{のろのろ},\ 
\text{ぶらぶら},\ \text{よたよた},\ \text{よろよろ}\}
\]
からなる．各サンプルは時間長 $T$ の 2 次元関節列として
\[
X \in \mathbb{R}^{T \times 14 \times 2}
\]
の形で与えられる．

\paragraph{座標の中心化とリサンプリング}
空間的な平行移動成分を除去するため，各時刻 $t$ ごとに 14 関節の重心
$\mathrm{COM}_t \in \mathbb{R}^{1 \times 1 \times 2}$ を求め，
\[
\tilde{X}_t = X_t - \mathrm{COM}_t
\]
として中心化された相対座標列 $\tilde{X}$ を得る．
その後，時間方向を線形補間により $T^\ast = 60$ フレームにリサンプリングし，
最終的な入力を
\[
\tilde{X} \in \mathbb{R}^{T^\ast \times 14 \times 2}
\]
に統一する．

\paragraph{正規化}
全サンプル（train split）の結合テンソルに対して，
全時刻・全関節・全サンプルに渡る平均 $\mu \in \mathbb{R}^{2}$，
分散 $\sigma^2 \in \mathbb{R}^{2}$ を計算し，
各サンプルを
\[
\hat{X} = (\tilde{X} - \mu) / \sigma
\]
で正規化する．同じ $\mu, \sigma$ を test split にも適用し，
学習時と評価時の分布ミスマッチを避ける．

\paragraph{入力統計のサニティチェック}
簡易解析スクリプトにより，各オノマトペごとに座標の標準偏差と
平均速度%
（隣接フレーム差分のノルムの平均）を計測した．
標準偏差は全ラベルでほぼ同程度である一方，
平均速度は
「どっしどっし」「のしのし」「よろよろ」などで高く，
「通常」「すたすた」「とぼとぼ」などで低い傾向が見られた．
これは，速度や振幅といった物理量に基づく粗いスタイル差は存在するが，
11 クラスを完全に分離するには情報が限定的であることを示唆している．

\paragraph{データセット拡張の可能性：Human3.6M へのオノマトペキャプショニング}
HOYO データセットのみでは，coarse 4 スタイルの分離は可能であるものの，
より多様な動き（走る，ジャンプ，転倒など）や，
fine-grained なオノマトペの区別を学習するにはデータ量が不足する可能性がある．
その場合の拡張案として，Human3.6M~\cite{human36m} などの大規模モーションデータセットを
オノマトペラベルでアノテーションする方法を検討する．

具体的には，Human3.6M の各モーションクリップ（関節軌跡）に対して，
以下のいずれかの方法でオノマトペラベルを付与する：
\begin{enumerate}
  \item \textbf{LLM による自動キャプショニング}：
    関節軌跡の統計的特徴（平均速度，加速度，COM の揺れ幅など）を抽出し，
    GPT-4 や Claude などの大規模言語モデルに対して
    「この動きを日本語のオノマトペで表現すると？」というプロンプトを送り，
    生成されたオノマトペをラベルとして採用する．
    もしくは，既存の 11 オノマトペから最も近いものを選択するタスクとして定式化する．
  \item \textbf{手動アノテーション}：
    モーションクリップを可視化（関節アニメーション）して，
    人間のアノテーターがオノマトペを付与する．
    これは LLM の結果を検証・補正する用途でも有用である．
\end{enumerate}

Human3.6M の関節構造は HOYO と異なるため，
Human3.6M の 3D 関節位置を HOYO 形式（2D，14 関節）に
リターゲットする必要がある．
具体的には，Human3.6M の主要関節（腰，股関節，膝，足首など）を
HOYO の 14 関節に対応させ，3D 座標を適切な視点から 2D 投影する．
このリターゲット処理により，Human3.6M の多様な動きを
MotionCLIP エンコーダに入力可能な形式に変換できる．

この拡張により，HOYO に含まれない「走る」「跳ぶ」「転ぶ」などの
動きもオノマトペラベルと共に学習データに追加でき，
MotionCLIP latent 空間の表現力と汎化性能が向上することが期待される．
実装の優先順位としては，まず HOYO のみでの学習結果を評価し，
coarse 4 スタイルの分離が不十分，または Navila 統合時に
スタイル報酬が機能しない場合に，Human3.6M 拡張を検討する．

\subsection{モーションエンコーダ：MotionCLIP VAE の活用}

モーション側の表現には，MotionCLIP~\cite{motionclip} の
トランスフォーマベースのオートエンコーダを用いる．
ただし本研究では，MotionCLIP 本来の CLIP 空間との整合性を維持するのではなく，
「事前学習済みの強力なモーション VAE」として利用する．

\paragraph{構成}
MotionCLIP のエンコーダ $E$ とデコーダ $D$ をベースに，
入力・出力次元を HOYO に合わせて再構成する．
具体的には，
関節数 $J=14$，特徴数 $C=2$，フレーム数 $T^\ast=60$ とし，
姿勢表現は SMPL ではなく 2 次元座標（$xyz$ に相当）とする．
エンコーダは $\hat{X}$ を受け取り，
潜在表現
\[
z_m = E(\hat{X}) \in \mathbb{R}^{D},\quad D=512
\]
を出力し，デコーダはこれを用いて
\[
\widehat{X} = D(z_m) \in \mathbb{R}^{T^\ast \times 14 \times 2}
\]
を再構成する．

MotionCLIP 内部の損失（関節姿勢，速度，頂点位置の再構成）を
まとめて $\mathcal{L}_{\mathrm{VAE}}$ とし，
既存の重みを初期値として HOYO データに対して微調整を行う．

\subsection{意味エンコーダ：オノマトペのテキスト埋め込み}

オノマトペの意味表現には 2 種類のテキスト埋め込みを用いる．

\paragraph{Sarashina 文埋め込み}
Sarashina モデル \texttt{sbintuitions/sarashina-embedding-v2-1b} を用い，
各ラベル $w$ に対して
\[
t(w) =
\begin{cases}
\text{``普通に歩いている。''} & (w = \text{通常}) \\
\text{``}w\text{と歩いている。''} & (\text{それ以外})
\end{cases}
\]
という日本語文を生成し，Sentence Transformers により
\[
z_{\mathrm{sem}}^{(\mathrm{fine})}(w) \in \mathbb{R}^{D_{\mathrm{sem}}}
\]
を得る．ベクトルは L2 正規化して用いる．

\paragraph{SigLIP テキストエンコーダ}
多言語 SigLIP モデル \texttt{google/siglip-base-patch16-256-multilingual} の
テキストエンコーダ（\texttt{SiglipTextModel}）を用い，同様のテキスト
$t(w)$ をトークナイズして得られる系列から
\[
z_{\mathrm{sem}}^{(\mathrm{fine})}(w)
= \mathrm{SigLIP\_Text}(t(w))_{[\mathrm{CLS}]}
\]
を抽出し，L2 正規化する．
これにより，視覚タスクに整合的な意味空間を利用できる．

\paragraph{粗いスタイルラベルへの集約}
11 語は，以下の 4 つのスタイル群にマージする：
\begin{itemize}
  \item 速い系：すたすた，せかせか，てくてく
  \item 遅い系：とぼとぼ，のろのろ
  \item 重い系：どっしどっし，のしのし
  \item ふらふら系：ぶらぶら，よたよた，よろよろ
\end{itemize}
各スタイル群 $g$ の意味ベクトルは，
対応する fine ラベルの平均として
\[
z_{\mathrm{sem}}^{(\mathrm{coarse})}(g)
= \frac{1}{|S_g|} \sum_{w \in S_g}
z_{\mathrm{sem}}^{(\mathrm{fine})}(w)
\]
と定義する．

\subsection{意味$\to$モーション射影と対照学習}

\paragraph{意味射影層}
テキスト側の意味空間と MotionCLIP latent を結ぶため，
線形射影
\[
f_{\mathrm{sem}}: \mathbb{R}^{D_{\mathrm{sem}}} \to \mathbb{R}^{D}
\]
を導入する．実装上はバイアスなし線形層
$\mathrm{sem\_proj} \in \mathbb{R}^{D \times D_{\mathrm{sem}}}$ であり，
各ラベルの意味ベクトルは
\[
z_s^{(k)} = \mathrm{normalize}
\bigl(
 \mathrm{sem\_proj}(z_{\mathrm{sem}}^{(k)})
\bigr)
\in \mathbb{R}^{D}
\]
として MotionCLIP latent 空間上のクラスプロトタイプを形成する．

\paragraph{教師ありコントラスト損失}
バッチ内のサンプル $(\hat{X}_i, y_i)$ に対して，
MotionCLIP エンコーダから得られる潜在表現
\[
z_m^{(i)} = \mathrm{normalize}(E(\hat{X}_i)) \in \mathbb{R}^{D}
\]
と，対応する意味プロトタイプ
\[
z_s^{(y_i)} \in \mathbb{R}^{D}
\]
を用い，教師ありコントラスト損失を定義する．

具体的には，
モーションと意味のペアを連結して
\[
\mathbf{f} =
\begin{bmatrix}
z_m^{(1)} \\ \vdots \\ z_m^{(B)} \\
z_s^{(y_1)} \\ \vdots \\ z_s^{(y_B)}
\end{bmatrix}
\in \mathbb{R}^{2B \times D},
\quad
\mathbf{y}_{\mathrm{ext}} =
\begin{bmatrix}
y_1 \\ \vdots \\ y_B \\
y_1 \\ \vdots \\ y_B
\end{bmatrix}
\]
とし，類似度行列
\[
S = \mathbf{f} \mathbf{f}^\top \in \mathbb{R}^{2B \times 2B}
\]
を計算する．温度パラメータ $T = \exp(\alpha)$
（$\alpha$ は学習可能な \texttt{logit\_scale}）を用いて
\[
\tilde{S}_{ij} = \exp(\alpha)\, S_{ij}
\]
とし，自己類似度（$i=j$）を除外した上で InfoNCE 形式の損失を構成する．
ラベルが等しいペアを正例，それ以外を負例とみなし，各 $i$ について
\[
p_i = 
\frac{\sum_j \exp(\tilde{S}_{ij}) \cdot \mathbf{1}[y_i = y_j,\ i\neq j]}
     {\sum_{j\neq i} \exp(\tilde{S}_{ij})}
\]
としたとき，
\[
\mathcal{L}_{\mathrm{contrast}}
= -\frac{1}{|\mathcal{I}|} \sum_{i \in \mathcal{I}} \log(p_i)
\]
をコントラスト損失とする
（有効な正例が存在する $i$ の集合を $\mathcal{I}$ とする）．

\paragraph{最終損失}
最終的な学習目標は
\[
\mathcal{L}
= \lambda_{\mathrm{VAE}}\,\mathcal{L}_{\mathrm{VAE}}
+ \lambda_{\mathrm{contrast}}\,\mathcal{L}_{\mathrm{contrast}}
\]
であり，実験では $\lambda_{\mathrm{VAE}}=1.0$，
$\lambda_{\mathrm{contrast}}\in\{0.3, 0.5\}$ を主に用いた．

\subsection{学習スケジュール}

\paragraph{freeze 段階}
初期段階では MotionCLIP のエンコーダ・デコーダを凍結し，
$\mathrm{sem\_proj}$ と $\alpha$ のみを学習する．
これにより，MotionCLIP の latent 構造を保ったまま，
テキスト埋め込み空間を適切に射影する初期解を得る．

\paragraph{encoder 段階}
次に，\texttt{stage=encoder} として
MotionCLIP エンコーダも学習対象に含める．
freeze 段階で得た
$\mathrm{sem\_proj}$ と $\alpha$ の最良モデルを初期値としてロードし，
HOYO データとオノマトペラベルに対して latent 空間を微調整する．

\subsection{評価と分析}

\paragraph{分類精度と再構成誤差}
評価関数 \texttt{evaluate\_dataset} は，
データセット全体を用いて
Top-1/Top-3 精度と MPJPE 風の再構成誤差を算出する．
各サンプル $i$ に対して
\[
\hat{y}_i = \arg\max_k \langle z_m^{(i)}, z_s^{(k)} \rangle
\]
により予測ラベルを得る．

fine 11 クラスの場合，全体の Top-1 精度は約 0.15 に留まり，
特に「よろよろ」と「のろのろ」を除く 9 クラスの精度はほぼ 0 であった．
一方，coarse 4 スタイルの場合，
Sarashina 埋め込みでは
Overall Acc@1 $\approx 0.46$，
各クラスの Acc@1 は
速い系 $\approx 0.51$，遅い系 $\approx 0.49$，
重い系 $\approx 0.48$，ふらふら系 $\approx 0.35$
と，粗いスタイルレベルでは有意な分離が得られた．

SigLIP 埋め込みを用いた場合，
Overall Acc@1 は $\approx 0.34$ であるものの，
ふらふら系の精度（$\approx 0.56$）は向上し，
代わりに遅い系の精度が低下するなど，
スタイルごとの得意・不得意が異なることが確認された．

\paragraph{潜在空間の可視化}
PCA による 2 次元可視化では，
coarse 4 スタイルにおいて，
速い系・遅い系・重い系・ふらふら系が
それぞれ緩やかなクラスタを形成し，
特に極端なスタイル（重い系や強いふらつき）に対応するサンプルが
特徴的な領域に集まる様子が観察された．

これらの結果から，
HOYO データのスケールと曖昧さを考えると
11 クラスの完全分類は困難である一方，
MotionCLIP latent 上で 4 種類の粗いスタイルを
区別することは十分可能であり，
強化学習におけるスタイル報酬
$\cos(z_{\mathrm{mot(agent)}}, z_{\mathrm{sem(onm)}})$
の基盤としては有用な表現が得られていると考えられる．

\section{Navila との統合に向けた具体的設計}

本節では，「将来的な構想」ではなく，
本研究の次フェーズとして必ず実施する Navila 連携実験について，
実装レベルで必要な要素と手順を整理する．

\subsection{環境設定とスタイル条件付き方策}

まず，強化学習の環境として IsaacLab 上の
H1 ヒューマノイドタスク \texttt{h1\_matterport\_vision} を用いる．
観測 $s_t$ は，Navila~\cite{navila} に倣い，
RGB-D 画像から抽出される視覚特徴ベクトルと，
関節角度・角速度などのロボット状態を結合したベクトルとする．
方策は，スタイル latent を条件として受け取る
\[
  \pi_\theta(a_t \mid s_t, z_{\mathrm{onm}})
\]
という形のニューラルネットワークとする．

ここで $z_{\mathrm{onm}}$ は，
Sarashina もしくは SigLIP による日本語テキスト埋め込み
$e_{\mathrm{text}}$ を，joint 学習で獲得した
線形射影 $\mathrm{sem\_proj}$ で MotionCLIP 空間に写像したもの
\[
  z_{\mathrm{onm}}
  = \mathrm{sem\_proj} \bigl( e_{\mathrm{text}}(w) \bigr)
\]
として実装する．
実際のネットワーク実装では，
$z_{\mathrm{onm}}$ を $s_t$ 側の特徴と連結するか，
FiLM のような条件付き正規化で中間特徴に注入する．

エピソード開始時に，タスク指示と一緒に
「すたすた」「のろのろ」といったオノマトペ $w$ をサンプリングし，
これを固定のスタイル条件 $z_{\mathrm{onm}}$ として
エピソード全体に与えることで，
「同じゴール地点でも質感の異なる歩き方」を学習させる．

\subsection{スタイル報酬の定義と計算手順}

Navila のタスク報酬 $r_{\mathrm{task},t}$ に加え，
MotionCLIP latent に基づくスタイル報酬を逐次計算して付与する．

\paragraph{H1 関節軌跡の HOYO 形式へのリターゲット（Option A）}
H1 ヒューマノイドの roll-out から得られる関節軌跡を，
MotionCLIP エンコーダに入力可能な HOYO 形式（2D，14 関節，60 フレーム）に変換する．
このリターゲット処理は以下の手順で行う：

\begin{enumerate}
  \item \textbf{座標系の正規化}：
    H1 の 3D 関節位置を，pelvis を原点とし，
    前方方向を $+x$ 軸，床面を $x$--$z$ 平面とする座標系に変換する．
    その後，$y$ 軸（高さ方向）を無視して $x$--$z$ 平面に投影し，
    2D 座標 $(x, z)$ を得る．
  \item \textbf{関節対応の定義}：
    H1 の主要関節を HOYO の 14 関節に対応させる．
    歩行スタイルに重要な下半身を優先し，
    腰（pelvis），両股関節（hip），両膝（knee），両足首（ankle）を
    HOYO の対応する関節にマッピングする．
    上半身（腕など）は初期実装では無視するか，
    簡単な対応（例：肩→腰の相対位置で近似）で処理する．
  \item \textbf{時間窓の抽出とリサンプリング}：
    直近 60 ステップの関節軌跡を抽出し，
    必要に応じて線形補間により HOYO の 60 フレーム形式に統一する．
    これにより，Isaac Sim の制御周期（約 50 Hz）と
    HOYO のフレームレート（約 30 Hz）の違いを吸収する．
  \item \textbf{正規化の適用}：
    HOYO データの学習時に用いた平均 $\mu$ と標準偏差 $\sigma$ を
    同じく適用して正規化し，最終的な入力
    $\hat{X} \in \mathbb{R}^{60 \times 14 \times 2}$ を得る．
\end{enumerate}

このリターゲット処理により，H1 の 3D 関節軌跡を
HOYO 形式に変換できるが，3D→2D 投影により情報が失われる点，
および H1 の腕の動きが反映されない点は注意が必要である．
将来的には，Option C（latent 空間での body-agnostic 化）により
これらの制約を緩和することを検討する．

\paragraph{スタイル報酬の計算}
時間窓 $[t, t+T)$ の関節軌跡 $x_{t:t+T}$ を
上記のリターゲット処理により HOYO 形式 $\hat{X}_{t:t+T}$ に変換したうえで
MotionCLIP エンコーダに入力し，エージェントの運動 latent
\[
  z_{\mathrm{mot(agent)}, t}
  = f_{\mathrm{mc\_enc}}(\hat{X}_{t:t+T})
\]
を計算する．
これと指示オノマトペの意味 latent $z_{\mathrm{sem(onm)}}$ との
コサイン類似度に基づき，
\[
  r_{\mathrm{style}, t}
  = \beta
    \cos\left(
      z_{\mathrm{mot(agent)}, t},\ 
      z_{\mathrm{sem(onm)}}
    \right)
\]
を定義する．ここで $\beta$ はタスク報酬との相対重みである．

実装上は，\textbf{毎ステップ sliding window 方式}で $r_{\mathrm{style}, t}$ を更新する．
具体的には，リングバッファに直近 60 フレームの関節軌跡を保持し，
各ステップ $t$ で以下を実行する：
\begin{enumerate}
  \item 現在の関節状態をリターゲットして HOYO 形式に変換し，バッファに追加する．
  \item バッファが 60 フレームに達したら，
        その内容を MotionCLIP エンコーダに入力して
        $z_{\mathrm{mot(agent)}, t}$ を計算し，
        $r_{\mathrm{style}, t} = \beta \cos(z_{\mathrm{mot(agent)}, t}, z_{\mathrm{sem(onm)}})$ を更新する．
  \item バッファが未満の場合は $r_{\mathrm{style}, t} = 0$ とする
        （もしくは，zero padding して早めに MotionCLIP を実行しても良い）．
\end{enumerate}

この方式により，スタイル報酬の応答性を保ちつつ，
MotionCLIP の入力形式（60 フレーム固定）を維持できる．
最終的な一段時間あたりの報酬は
\[
  r_t = r_{\mathrm{task}, t} + r_{\mathrm{style}, t} + r_{\mathrm{reg}, t}
\]
とし，$r_{\mathrm{reg}, t}$ には姿勢崩壊や転倒を抑制する正則化項を含める．

\subsection{PPO と対照学習の joint 更新スケジュール}

学習ループは，PPO と対照学習を同一ロールアウトから更新する形で構成する．
具体的な一反復の流れは次の通りである．
\begin{enumerate}
  \item Navila の高レベルプランナから
        ゴール位置とオノマトペ $w$ を受け取り，
        $z_{\mathrm{onm}}$ を計算して環境にリセットする．
  \item Policy $\pi_\theta(a_t \mid s_t, z_{\mathrm{onm}})$ を用いて
        $N$ ステップのロールアウトを収集し，
        各ステップで $r_{\mathrm{task}, t}$ と
        $r_{\mathrm{style}, t}$（sliding window 方式で毎ステップ更新）を計算する．
  \item ロールアウトから PPO の損失 $\mathcal{L}_{\mathrm{PPO}}(\theta)$ を計算し，
        $\theta$ を更新する．
  \item \textbf{対照学習の更新戦略（Phase 1：HOYO のみ固定）}：
    初期実装（Phase 1）では，MotionCLIP エンコーダと $\mathrm{sem\_proj}$ は
    \textbf{完全に freeze のまま}とし，RL 中には更新しない．
    これにより，HOYO データで学習済みの latent 空間を保ったまま，
    Option A リターゲットによる $r_{\mathrm{style}}$ のみで Navila を学習する．
    
    将来的な Phase 2 では，オンライン roll-out データを
    対照学習に混ぜることを検討するが，その際は以下の注意点がある：
    \begin{itemize}
      \item 学習初期のオンライン軌跡は「転倒」「停止」などが多く，
            HOYO の歩行データと分布が大きく異なるため，
            これを混ぜると MotionCLIP latent が壊れる可能性がある．
      \item 対策として，まずはタスク成功率がある閾値（例：50--60\%）を
            超えたエピソードのみを「良質サンプル」として採用する．
      \item オンライン比率を段階的に増やす（$0\% \rightarrow 10\% \rightarrow 30\% \rightarrow 50\%$）．
      \item 更新対象は，まず $\mathrm{sem\_proj}$ のみとし，
            問題なければ MotionCLIP エンコーダの最後の数層のみを解凍する．
    \end{itemize}
    
    Phase 1 では，対照学習の更新ステップ（上記の 4 番目のステップ）は
    実行せず，PPO のみで学習を進める．
    これにより，「HOYO で学習済み MotionCLIP を固定して，
    Option A リターゲットだけで Navila を動かしたらどうなるか」を
    まず検証する．
\end{enumerate}

全体として最小化する目的関数は
\[
  \mathcal{L}_\mathrm{total}
  = \mathcal{L}_{\mathrm{PPO}}(\theta)
  + \lambda\,\mathcal{L}_{\mathrm{contrast}}
\]
とし，$\lambda$ はスタイル整合性をどの程度重視するかを制御するハイパパラメータである．
ただし Phase 1 では $\lambda = 0$（すなわち $\mathcal{L}_{\mathrm{contrast}}$ は計算しない）とし，
MotionCLIP は完全 freeze のまま PPO のみで学習する．

Phase 2 以降では，MotionCLIP 側の更新戦略として，
\texttt{stage=freeze}（RL から参照するだけ），
\texttt{stage=encoder}（エンコーダのみ微調整），
\texttt{stage=full}（エンコーダ＋デコーダを RL データで追加学習）
の三段階を順に採用し，
各段階で Navila の成功率とスタイル報酬のバランスを比較する．

\subsection{Navila パイプラインへの組み込みと評価プロトコル}

Navila~\cite{navila} 全体のパイプラインにおいては，
高レベルの言語指示から
\begin{enumerate}
  \item 目的地に関するテキスト（例：``リビングのソファへ行って''）と，
  \item 歩行質感に関するオノマトペ（例：``すたすた''） 
\end{enumerate}
の二種類を VLM に入力し，
それぞれからゴール位置と $z_{\mathrm{onm}}$ を生成する構成を取る．
前者は既存の Navila の枠組みを踏襲し，
後者のみ本研究の MotionCLIP ベーススタイルモジュールで処理する．

評価指標としては，
\begin{itemize}
  \item タスク成功率（目標地点への到達率），
  \item 平均エピソード長や衝突回数などの従来指標，
  \item MotionCLIP latent 空間上での
        $\cos(z_{\mathrm{mot(agent)}}, z_{\mathrm{sem(onm)}})$ の分布，
  \item coarse 4 スタイルラベルに対する混同行列，
        および latent PCA 上でのクラスタ分離度
\end{itemize}
を用いる．
これにより，
「ゴールへたどり着けているか」と
「指示した質感で歩けているか」を
同時に定量評価できるようにする．

実装タスクとしては，
MotionCLIP joint モデルの IsaacLab/NaVILA へのデプロイ，
スタイル報酬計算モジュールの実装，
joint 更新ループの構築，
および上記指標に基づく評価スクリプトの整備を行えば，
本節で述べた一連の質感ナビゲーション実験をそのまま実行できる状態となる．

\paragraph{実装マイルストーンとデータ拡張の判断基準}
実装は以下の段階で進める：
\begin{enumerate}
  \item \textbf{Week 1--2：MotionCLIP 単体の最終検証}
    HOYO データのみで学習した joint モデル（coarse 4 スタイル）の
    潜在空間を PCA/t-SNE で可視化し，
    クラスタ分離度を定量的に評価する．
    もし coarse 4 スタイルの分離が不十分（例：混同行列の対角成分が 0.4 未満），
    または Navila 統合時にスタイル報酬が機能しない場合，
    Human3.6M へのオノマトペキャプショニングによるデータ拡張を検討する．
  \item \textbf{Week 3--4：Isaac Sim 環境構築と Option A retargeting}
    H1 Matterport タスクを動かし，
    H1 の関節軌跡を HOYO 形式（2D，14 関節）にリターゲットするモジュールを実装する．
    手動で「速く歩く」「ゆっくり歩く」を実行し，
    その latent が HOYO データのクラスタに適切に落ちるかを確認する．
  \item \textbf{Week 5--6：スタイル報酬モジュールの実装}
    60 フレーム sliding window 方式で $r_{\mathrm{style}}$ を計算するモジュールを実装し，
    ランダムポリシーや簡単な PD 歩行ポリシーで動作確認する．
    $\beta \in \{0, 0.1, 0.5, 1.0\}$ の ablation で報酬スケールを検証する．
  \item \textbf{Week 7--8：PPO 学習（MotionCLIP freeze）}
    Navila ベースの PPO に $r_{\mathrm{style}}$ を統合し，
    MotionCLIP は完全 freeze のまま学習する．
    タスク成功率とスタイル報酬の推移を記録し，
    指示したオノマトペに対して latent 類似度が向上するかを確認する．
  \item \textbf{Week 9--10：encoder 解凍と optional online contrastive}
    時間と成果次第で，MotionCLIP エンコーダの一部を解凍して
    RL 中に微調整する．
    オンラインデータを contrastive loss に混ぜる場合は，
    成功エピソードのみを採用し，オンライン比率を段階的に増やす．
\end{enumerate}

このマイルストーンに基づき，
Week 1--2 の段階で HOYO データのみでは不十分と判断した場合，
Human3.6M 拡張を Week 3--4 の並行タスクとして進める．
Human3.6M のリターゲット処理と LLM キャプショニングは
HOYO のみでの学習と並行して実装可能であるため，
早期に判断できれば遅延なく拡張データを投入できる．

\subsection{今後の拡張：音韻エンコーダと未知オノマトペ}

project 計画では，最終的に
phon/sem/motion の三本エンコーダを構成し，
音韻から意味 latent を推定する \textit{phon$\to$sem} マッピングを介して，
未知オノマトペに対してもモーション latent を生成することを構想している．
本稿で扱った sem$\leftrightarrow$motion の対照学習は，
この三本構造のうち sem/motion 間を先に整備した段階に相当する．
今後は，音韻表現（仮名列や音素系列）に対するエンコーダと，
phon latent から sem latent を再構成する MLP を導入し，
未知語に対しても
「音韻 $\to$ 意味 $\to$ モーション latent」
の経路を通じたスタイル付与を可能にする予定である．

さらに，MotionCLIP 論文が行っているように，
CLIP テキスト／画像エンコーダの latent 空間との alignment を
補助的な損失として取り込むことで，
より抽象的な言語指示や視覚的スタイル参照に対しても
ロバストな質感ナビゲーションが可能かどうかを検証したい．
\begin{thebibliography}{99}
  {\footnotesize
  \bibitem{navila}
  A.Cheng et al., ``NaVILA: Legged Robot Vision--Language--Action Model for Navigation,'' arXiv:2412.04453, 2025.\@

  \bibitem{motionclip}
  G.~Tevet et al., ``MotionCLIP: Exposing Human Motion Generation to CLIP Space,'' ECCV, 2022.\@
  }
\end{thebibliography}

  

\end{document}